\documentclass[11pt]{article}
\usepackage[margin=1.1in]{geometry}
\usepackage{amsmath,amssymb,amsthm}
\usepackage{booktabs,tabularx,array}
\usepackage{microtype}
\usepackage{xcolor}
\usepackage[colorlinks=true,linkcolor=blue!50!black,citecolor=blue!50!black,urlcolor=blue!50!black]{hyperref}
\newtheorem{theorem}{Theorem}[section]
\newtheorem{lemma}[theorem]{Lemma}
\newtheorem{proposition}[theorem]{Proposition}
\newtheorem{corollary}[theorem]{Corollary}
\theoremstyle{remark}
\newtheorem{remark}[theorem]{Remark}
\newcommand{\N}{\mathbb N}
\newcommand{\Z}{\mathbb Z}
\newcommand{\R}{\mathbb R}
\newcommand{\fr}[1]{\{#1\}}
\newcommand{\code}[1]{{\small\ttfamily #1}}
\setlength{\emergencystretch}{1.5em}
\title{An Unconditional Sturmian Exclusion\\for the Collatz Map at Slope $1/\sqrt2$}
\author{M.~Sharpe\thanks{The mathematics, Lean formalization, and exposition were developed with Codex (OpenAI), building on the author's existing Collatz repository. The Collatz conjecture remains open. No claim of priority in the mathematical literature is made.}}
\date{September 4, 2026}
\begin{document}
\maketitle
\begin{abstract}
Let $T(n)=n/2$ for even $n$ and $T(n)=(3n+1)/2$ for odd $n$.
We prove that no natural number has the mechanical parity itinerary
$\lfloor(t+1)\alpha+\rho\rfloor-\lfloor t\alpha+\rho\rfloor$ with
$\alpha=1/\sqrt2$, for any real intercept $\rho$. The same exclusion
holds after every finite orbit prefix. This is an unconditional instance
of the repository's previously conditional Sturmian level theorem.
A Pell recurrence supplies determinant-one bracketing pairs; an elementary
rational-grid argument supplies uniform visits to the relevant rotation
interval, using a slightly larger grid than the one suggested by the
three-distance theorem. The remaining size comparison reduces to a linear
factor times a geometric sequence tending to zero. All approximation,
coverage, and size hypotheses are discharged in Lean~4. The proof depends
only on Lean's standard foundational axioms and contains no \code{sorry},
\code{native\_\allowbreak{}decide}, or cited mathematical axiom. We also record the
finite-word and survivor-cylinder identities formalized alongside this
result, distinguishing these tools from any unproved universal coverage
principle.
\end{abstract}

\section{Statement and scope}
We use the shortcut, or Terras, map
\[
T(n)=\begin{cases}n/2,&n\equiv0\pmod2,\\(3n+1)/2,&n\equiv1\pmod2,\end{cases}
\qquad n\in\N=\{0,1,2,\ldots\}.
\]
For positive starts it compresses the mandatory halving after an odd step
of the ordinary Collatz map. Write $n_t=T^t(n)$ and $w_t(n)=n_t\bmod2$.
For real $\alpha,\rho$, define the lower mechanical word
\begin{equation}\label{eq:mechanical}
w_{\alpha,\rho}(t)=\lfloor(t+1)\alpha+\rho\rfloor-\lfloor t\alpha+\rho\rfloor.
\end{equation}
When $0<\alpha<1$, its letters are zero or one. Telescoping gives
$\sum_{t<s}w_{\alpha,\rho}(t)=\lfloor s\alpha+\rho\rfloor-\lfloor\rho\rfloor$;
in particular its limiting density of ones is $\alpha$.

\begin{theorem}\label{thm:main}
Set $\alpha=\sqrt2/2$. For every $N\in\N$ and $\rho\in\R$, there exists
$t\in\N$ such that $T^t(N)\bmod2\ne w_{\alpha,\rho}(t)$.
\end{theorem}
\begin{corollary}\label{cor:tails}
For every $N,s\in\N$ and $\rho\in\R$, the equality
$T^{s+t}(N)\bmod2=w_{\alpha,\rho}(t)$ cannot hold for all $t\in\N$.
\end{corollary}
\begin{proof}
Apply Theorem~\ref{thm:main} to the natural number $T^s(N)$.
\end{proof}
The theorem is \code{Collatz.Silver.no\_\allowbreak{}itinerary}; the corollary is
\code{Collatz.Silver.no\_\allowbreak{}eventual\_\allowbreak{}itinerary}. The formal slope is
\code{Real.sqrt 2 / 2}. No assumption of divergence is needed in either
statement. The inequality $\alpha>2/3>\log 2/\log 3$ puts this family above
the critical density, so the result is not a subcritical drift argument.

The previous work \cite{rung1} supplied a conditional Sturmian level
theorem and a proposed continued-fraction route to its hypotheses.
Here we prove a specific unconditional case and two reusable arithmetic
inputs. We do not exclude every Sturmian slope, classify every possible
counterexample itinerary, or eliminate all nontrivial cycles.

\section{The orbit-theoretic reduction}
We recall the formal ingredients used below, including the correction
term that must be retained when estimating a shifted orbit.
Let $j_s(n)=\sum_{t<s}w_t(n)$. There is a nonnegative integer $d_s(n)$ with
\begin{equation}\label{eq:affine}
2^sT^s(n)=3^{j_s(n)}n+d_s(n),\qquad
 d_s(n)=\sum_{i<s}w_i(n)\,2^i3^{j_s(n)-j_{i+1}(n)}.
\end{equation}
These are \code{terras\_\allowbreak{}exact\_\allowbreak{}form} and \code{dcoef\_\allowbreak{}closed}
\cite{shadow,ladder}. The parity-vector bijection states that two starts
have the same first $m$ parities if and only if they are congruent modulo
$2^m$. Its two directions follow inductively from the doubling identities
for $T$ and the invertibility of $3$ modulo powers of $2$.

\begin{lemma}[Prefix-power bound]\label{lem:power}
Suppose the first $M$ parity letters of $n$ are $\ell$-periodic, where
$M\ge\ell$. If $T^\ell(n)\ne n$, then
\begin{equation}\label{eq:power}
2^M\le 2^\ell n+3^o(n+2^\ell),\qquad o=j_\ell(n).
\end{equation}
\end{lemma}
\begin{proof}
The starts $n$ and $T^\ell(n)$ share their first $M-\ell$ parities.
Their nonzero difference is divisible by $2^{M-\ell}$, so
$2^{M-\ell}\le\max(n,T^\ell(n))$. The established integer growth bound
$2^\ell T^\ell(n)\le3^o(n+2^\ell)$, together with
$\max(n,T^\ell(n))\le n+T^\ell(n)$, gives \eqref{eq:power}.
This is \code{prefix\_\allowbreak{}power\_\allowbreak{}bound} in \cite{rung1}.
\end{proof}

\begin{lemma}[Mechanical-word tail bound]\label{lem:tail}
Suppose $N$ has itinerary \eqref{eq:mechanical}, with $0\le\alpha$ and
$2\le3^\alpha$. Then
\begin{equation}\label{eq:tail}
2^sT^s(N)\le3^{s\alpha}(3N+3s)\qquad(s\in\N).
\end{equation}
\end{lemma}
\begin{proof}
Every window of length $r$ contains at most $r\alpha+1$ ones. Thus
$3^{j_s(N)}\le3\cdot3^{s\alpha}$. In the sum in
\eqref{eq:affine}, each term is at most $3\cdot3^{s\alpha}$: use
$2^i\le3^{i\alpha}$ and the window bound on ones following position $i$.
Hence $d_s(N)\le3s\cdot3^{s\alpha}$, proving \eqref{eq:tail}.
The formal declaration is \code{tail\_\allowbreak{}le\_\allowbreak{}of\_\allowbreak{}itin}.
\end{proof}

For an approximation $q\alpha=p+\delta$, where $q\in\N$, $p\in\Z$,
and $0<\delta<1$, call $n$ a \emph{visit} if
\[
\fr{n\alpha+\rho}\ge1-\delta.
\]
Here $\fr{x}=x-\lfloor x\rfloor$ is the fractional part. If $V(n)$ is the
indicator of a visit, subtraction of the two floor identities gives
\begin{equation}\label{eq:shift}
w_{\alpha,\rho}(n+q)-w_{\alpha,\rho}(n)=V(n+1)-V(n).
\end{equation}

\begin{proposition}[Level theorem with an endpoint test]\label{prop:level}
Suppose $0\le\alpha$, $2\le3^\alpha$, $q\alpha=p+\delta$ with
$0<\delta<1$, and $Q\ge3$. Assume
\begin{align}
&|m\alpha-r|\ge\delta
 &&(0<m<Q,\ r\in\Z),\label{eq:sep}\\
&\text{every window }[a,a+G)\text{ contains a visit},\label{eq:vis}\\
&3^{(q+G)\alpha+1}(3N+3G+1)<2^{Q-3+G}.\label{eq:endpoint}
\end{align}
Then $N$ does not have itinerary $w_{\alpha,\rho}$.
\end{proposition}
\begin{proof}
Two visits are at least $Q$ apart: otherwise their fractional parts,
both in the half-open interval $[1-\delta,1)$, differ by less than
$\delta$, contradicting \eqref{eq:sep}. Let $n_0<G$ be the first visit
and put $s=n_0+1\le G$. There are no visits strictly between $n_0$ and
$n_0+Q$. Equation~\eqref{eq:shift} therefore makes the tail beginning at
$s$ periodic with period $q$ on its first $q+Q-2$ letters.

This tail is not a $q$-periodic orbit. If it were, \eqref{eq:shift} would
make $V(n)$ constant for $n\ge s$. Coverage forces that constant to be
one; then adjacent visits contradict separation. Apply
Lemma~\ref{lem:power} to $m=T^s(N)$, using $o\le q\alpha+1$ and
$2^q\le3^{q\alpha}$. A convenient weakening is
\[
2^{q+Q-2}\le2^{q+1}3^{q\alpha+1}(m+1).
\]
After multiplying by $2^s$ and applying Lemma~\ref{lem:tail} and
$2^s\le3^{s\alpha}$, cancellation gives
\begin{equation}\label{eq:forced}
2^{Q-3+s}\le3^{(q+s)\alpha+1}(3N+3s+1).
\end{equation}
However the ratio of the right side to the left side in
\eqref{eq:forced} is nondecreasing in $s$: up to a positive constant it
is $(3^\alpha/2)^s(3N+3s+1)$. Thus \eqref{eq:endpoint} implies that this
ratio is less than one for every $s\le G$, a contradiction.
\end{proof}
The last monotonicity step is \code{sturmian\_\allowbreak{}size\_\allowbreak{}at\_\allowbreak{}endpoint}.
The proposition is \code{sturmian\_\allowbreak{}level\_\allowbreak{}endpoint}, built on the
existing \code{sturmian\_\allowbreak{}level}.

\section{Two elementary approximation lemmas}
The following inputs replace the general continued-fraction machinery for
our chosen recurrence. They are proved independently of the Collatz map.

\begin{lemma}[Separation from a bracketing pair]\label{lem:neighbors}
Let $q,Q>0$ be integers and $p,P\in\Z$. Suppose
\[
q\alpha=p+\delta,\qquad Q\alpha=P-\varepsilon,\qquad
\delta,\varepsilon>0,\qquad Pq-pQ=1.
\]
Then \eqref{eq:sep} holds.
\end{lemma}
\begin{proof}
Given $0<m<Q$ and $r\in\Z$, set
$A=Pm-Qr$ and $B=qr-pm$. The determinant identity gives
\[
m=Aq+BQ,\qquad r=Ap+BP,\qquad m\alpha-r=A\delta-B\varepsilon.
\]
If $A\le0$, positivity of $m$ forces $B\ge1$. Moreover $A=0$ would give
$m=BQ\ge Q$, so $A\le-1$ and the error is at most $-\delta$.
If $A>0$, then $A\ge1$, while $m<Q$ forces $B\le0$; the error is at
least $\delta$. These cases prove the assertion.
\end{proof}

\begin{lemma}[Visits from a rational grid]\label{lem:grid}
Let $G>0$ be an integer and $R,U,V\in\Z$. Suppose
\begin{equation}\label{eq:griddata}
UR+VG=1,\qquad G\alpha=R+e,\qquad e\ge0,
\qquad 0<\delta<1,\qquad 1+Ge\le G\delta.
\end{equation}
Then for every real $\rho$, each window $[a,a+G)$ contains a visit.
\end{lemma}
\begin{proof}
First take $a=0$ and let $z=\lceil G(1-\delta-\rho)\rceil$.
Bezout's identity gives an integer $0\le n<G$ and $h\in\Z$ with
$nR=z+hG$: for example, choose $n\equiv zU\pmod G$.
Set $y=n\alpha+\rho-h$. Then
\[
Gy=z+G\rho+ne.
\]
The ceiling bounds and $0\le ne\le Ge$ imply
\[
G(1-\delta)\le Gy<G(1-\delta)+1+Ge\le G.
\]
Consequently $1-\delta\le y<1$, so $y=\fr{n\alpha+\rho}$ is a visit.
For general $a$, apply this argument with intercept $a\alpha+\rho$ and
then translate the index by $a$.
\end{proof}
The formal declarations are \code{sturmian\_\allowbreak{}separation\_\allowbreak{}of\_\allowbreak{}neighbors}
and \code{rotation\_\allowbreak{}window\_\allowbreak{}of\_\allowbreak{}grid}. The grid lemma handles all real
intercepts and the half-open interval endpoints directly. A larger $G$
can make the width condition easier to prove, at the expense of the size
condition in Proposition~\ref{prop:level}.

\section{Pell certificates for $\alpha=\sqrt2/2$}
Define natural-number pairs by
\begin{equation}\label{eq:pellrec}
(q_0,p_0)=(3,2),\qquad
q_{k+1}=3q_k+4p_k,\qquad p_{k+1}=2q_k+3p_k.
\end{equation}
Induction proves
\begin{equation}\label{eq:pell}
q_k^2-2p_k^2=1,\quad 3\le q_k,\quad 1\le p_k,\quad
q_k\le2p_k,\quad p_k\le q_k,\quad k+3\le q_k.
\end{equation}
Fix $k$ and abbreviate $q=q_k$, $p=p_k$, $\delta=q\alpha-p$.
Since $\alpha^2=1/2$, the Pell identity gives
\begin{equation}\label{eq:deltaproduct}
\delta(q\alpha+p)=\frac12.
\end{equation}
Thus $0<\delta<1$. We shall use
\[
\frac23<\alpha<\frac57,
\]
which follows by squaring positive quantities and comparing with $1/2$.

\subsection{Separation}
Put $Q=q+2p$ and $P=q+p$. Direct calculation gives
\[
Pq-pQ=1,\qquad Q\alpha=P-\varepsilon,
\qquad \varepsilon=(2\alpha-1)\delta>0.
\]
Lemma~\ref{lem:neighbors} supplies the separation bound through $Q-1$.

\subsection{Coverage}
Use the next lower approximation as the grid:
\[
G=q_{k+1}=3q+4p=2Q+q,\qquad R=p_{k+1}=2q+3p.
\]
Its error and Bezout certificate are
\begin{equation}\label{eq:next}
e=G\alpha-R=(3-4\alpha)\delta>0,
\qquad -QR+PG=1.
\end{equation}
The required width condition follows from the exact identity
\begin{equation}\label{eq:width}
G(\delta-e)-1
=\delta\bigl((10\alpha-6)q+(16\alpha-10)p\bigr)\ge0.
\end{equation}
To verify it, substitute \eqref{eq:next} and replace $1$ using twice
\eqref{eq:deltaproduct}. The last inequality follows from
$\alpha>2/3$ and nonnegativity of $q,p,\delta$.
Lemma~\ref{lem:grid}, with $U=-Q$ and $V=P$, proves uniform coverage.

\begin{center}
\begin{tabular}{rrrrr}
\toprule
$k$ & $q_k$ & $p_k$ & $Q=q_k+2p_k$ & $G=q_{k+1}$\\
\midrule
0 & 3 & 2 & 7 & 17\\
1 & 17 & 12 & 41 & 99\\
2 & 99 & 70 & 239 & 577\\
3 & 577 & 408 & 1393 & 3363\\
\bottomrule
\end{tabular}
\end{center}
Unlike a finite verification of candidate starting values, each row
supplies arithmetic statements quantified over every intercept. The
recurrence and its invariants provide these certificates for all levels.

\section{The size estimate and completion of the proof}
Put $a=3^\alpha$. We have $a\ge2$. A useful strict rational upper bound is
\begin{equation}\label{eq:abounds}
a<\frac{20}{9},\qquad a^2\le8,\qquad a^6<128.
\end{equation}
Indeed $a^7=3^{7\alpha}<3^5=243<(20/9)^7$, and the remaining comparisons
are rational arithmetic. The formal proof uses these inequalities,
not floating-point logarithms.

From \eqref{eq:pell}, $Q\ge2q$ and $G\le7q$. Let $d=Q-2q\ge0$. Then
\begin{equation}\label{eq:exponents}
q+G=6q+2d,\qquad Q+G=7q+3d.
\end{equation}
For a fixed candidate start $N$, multiply \eqref{eq:endpoint} by $8$.
As $Q\ge3$, it becomes
\begin{equation}\label{eq:multiplied}
24a^{q+G}(3N+3G+1)<2^{Q+G}.
\end{equation}
Using \eqref{eq:abounds} and \eqref{eq:exponents}, its left side is at most
\[
24(3N+21q+1)(a^6)^q8^d,
\]
whereas its right side is $128^q8^d$. Thus it suffices that
\begin{equation}\label{eq:scalar}
24(3N+21q+1)\left(\frac{a^6}{128}\right)^q<1.
\end{equation}
Since $0<a^6/128<1$, a geometric decay dominates the linear factor.
For each $N$, inequality~\eqref{eq:scalar} holds for all sufficiently
large $q$. The lower bound $k+3\le q_k$ ensures a suitable level exists.

\begin{proof}[Proof of Theorem~\ref{thm:main}]
Fix $N$ and $\rho$. Choose a level large enough for \eqref{eq:scalar}.
The Pell certificates supply \eqref{eq:sep}, the rational grid supplies
\eqref{eq:vis}, and the preceding estimate supplies
\eqref{eq:endpoint}. Proposition~\ref{prop:level} gives the contradiction.
\end{proof}
The only limit theorem needed in the final step is Mathlib's proved
statement that $n^j c^n\to0$ for $|c|<1$, here for $j=0,1$.
The size estimate is independent of $\rho$, while the grid argument
covers every $\rho$.

\section{Formalization and verification}
The new proof is organized as follows. All paths in the table are
relative to \code{lean/Collatz/}.
\begin{center}
\small
\begin{tabularx}{\linewidth}{@{}lX@{}}
\toprule
Module & Mathematical responsibility\\
\midrule
\code{SturmianEndpoint.lean} & Reduce all size inequalities to the endpoint.\\
\code{SturmianApprox.lean} & Determinant-one separation and Bezout-grid coverage.\\
\code{SturmianSilver.lean} & Pell recurrence, slope and width bounds, eventual size estimate, final exclusion.\\
\bottomrule
\end{tabularx}
\end{center}
They use the existing parity, shadow, prefix-power, and Sturmian modules.
The main declaration has the type
\begin{quote}\small\ttfamily
Collatz.Silver.no\_\allowbreak{}itinerary (N : $\N$) ($\rho$ : $\R$) :\\
\hspace*{1em}$\neg$ Collatz.HasItin Collatz.Silver.alpha $\rho$ N
\end{quote}
Here \code{HasItin} means equality of every parity letter with
\eqref{eq:mechanical}; it does not include an additional hypothesis
about approximation, density, or divergence.

The selected modules were rebuilt with Lean
\code{v4.27.0-rc1}. Querying \code{\#print axioms} for both headline
declarations returned exactly \code{propext}, \code{Classical.choice},
and \code{Quot.sound}. These are the standard foundational axioms used
by Mathlib, not assumptions of unproved Collatz statements. The
reproducible audit script \code{analysis/audit\_\allowbreak{}theorems.py} records
printed types, dependencies, and project provenance. This verification
was not an independent kernel replay of every imported library or a
fresh rebuild of the large KL certificates. The finite experiments are
not premises of the theorem.

\section{Additional formal arithmetic tools}\label{sec:tools}
The same research pass produced the following reusable statements. They
are included to document the associated results without presenting them
as a proof of universal itinerary coverage.

\subsection{Adjacent swaps and compensated differences}
For a Boolean word $u$, let $|u|$ be its length, $o(u)$ its number of ones,
and $d(u)$ its affine correction. The concatenation law is
\[
d(uv)=3^{o(v)}d(u)+2^{|u|}d(v).
\]
Since $d(01)=2$ and $d(10)=1$, it implies the exact identity
\begin{equation}\label{eq:swap}
d(u01v)=d(u10v)+2^{|u|}3^{o(v)}.
\end{equation}
If equal-length, equal-weight words $u,v$ are realized from starts $x,y$,
with length $L$ and weight $o$, their endpoints satisfy
\[
2^L(y'-x')=3^o(y-x)+d(v)-d(u).
\]
When $y'\ne x'$, the compensated integer on the right is nonzero and
has absolute value at least $2^L$. These are
\code{WordAffine.adjacent\_\allowbreak{}swap}, \code{compare}, and \code{compare\_\allowbreak{}bound}.
An actual exclusion still requires an independent upper bound on the
weighted correction difference. A small number of mismatches alone
provides no such bound.

\subsection{Height-limited survivor intervals}
For $r<2^k$ and a nonnegative quotient $h$, every prefix endpoint is
\begin{equation}\label{eq:cylinder}
T^i(r+2^kh)=T^i(r)+3^{j_i(r)}2^{k-i}h\qquad(0\le i\le k).
\end{equation}
This follows by comparing \eqref{eq:affine} on congruent starts.
Consequently, the height condition $1\le r+2^kh\le H$ and all prefix
survival conditions $T^i(r+2^kh)\ge r+2^kh$ are affine inequalities in
the same integer $h$. Their intersection is a finite integer interval,
possibly empty. Membership and order convexity are formalized in
\code{Cylinder.mem\_\allowbreak{}parameters} and \code{parameters\_\allowbreak{}convex};
\eqref{eq:cylinder} is \code{prefix\_\allowbreak{}transport}.
This describes finite survivor sets exactly, but does not prove they
eventually disappear at every height.

\subsection{Initial-height bounds must include the correction}
For an unbounded orbit with an $\ell$-periodic block of $M$ letters
starting at time $s$, let $o=j_\ell(T^s(n))$. Multiplying the prefix-power
bound by $2^s$ and using \eqref{eq:affine} gives
\begin{equation}\label{eq:seedbound}
\begin{split}
2^{M+s}\le{}&(2^\ell+3^o)\bigl(3^{j_s(n)}n+d_s(n)\bigr)\\
&+3^o2^{\ell+s}.
\end{split}
\end{equation}
This is \code{prefix\_\allowbreak{}power\_\allowbreak{}initial\_\allowbreak{}bound}. It supplies exact finite
excluded-height certificates for the revised repetition experiment.
Dropping $d_s(n)$ would be unjustified: $T(3)=5>3\cdot3/2$ already
contradicts the corresponding drift-only upper bound.

\subsection{A precise bounded-correction hypothesis}
Define $c_t=d_t(n)/3^{j_t(n)}$. The formal theorem
\code{supercritical\_\allowbreak{}of\_\allowbreak{}geometric} states that if $C\ge0$, $0\le r<1$,
and
\[
2^t\le Cr^t3^{j_t(n)}\qquad\text{for every }t,
\]
then $c_t\le C/(1-r)$ for every $t$. The proof bounds each correction
increment by $Cr^t$ and sums the geometric series, expressed as a finite
induction in Lean. This is an explicit sufficient condition. The
critical density floor alone does not supply it for every hypothetical
divergent orbit.

\newpage
\section{What remains open}
The present theorem rules out one slope with every intercept, and the
arithmetic lemmas make further instances plausible targets. Other slopes
need their own certified approximation and size data. The swap identities
need a genuinely useful bound on accumulated weighted defects. Neither
these tools nor the exact survivor intervals show that every possible
counterexample itinerary belongs to an excluded family. Nontrivial cycles
also remain a separate issue. The formal result is the quantified family
exclusion in Theorem~\ref{thm:main}, with these limits kept explicit.

\begin{thebibliography}{9}
\raggedright
\bibitem{rung1} M.~Sharpe, \emph{Rung One of the Word-Complexity Ladder:
a Machine-Verified Exclusion of an Automatic Collatz Itinerary by
Mahler's Method in Degree One}, August 2026.
Repository file \code{paper/rung1.tex}, especially its prefix-power
criterion and conditional Sturmian level theorem.
\bibitem{shadow} M.~Sharpe, companion paper on the 3-adic shadow,
\code{paper/shadow.tex}; formal development \code{lean/Collatz/Shadow.lean}.
\bibitem{ladder} M.~Sharpe, companion paper on the word-complexity ladder,
\code{paper/ladder.tex}; formal development \code{lean/Collatz/Ladder.lean}.
\bibitem{source} M.~Sharpe, \emph{Collatz at the Critical Line}, source
repository, \url{https://github.com/msharpe248/collatz}.
For this result see \code{SturmianSilver.lean} and its supporting modules;
verification metadata are in \code{analysis/theorem\_\allowbreak{}audit.json}.
\end{thebibliography}
\end{document}
